%%%%%%%%%%  scan idx 41  |  folio 28  %%%%%%%%%%
b)
Lorsque le corps $k$ n'est pas algébriquement clos, il est plausible que $\underline{K}(G)$ est un $\lambda$-anneau spécial en caractéristique $0$, et de même pour $\underline{K}_r(G)$ en toutes caractéristiques.

c) En caractéristique non nulle, le théorème 1.1 m'a été gracieusement fourni par Chevalley.

\subsection*{\S 3. Le $\lambda$-anneau défini par un anneau gradué}

Soit $A$ un anneau gradué commutatif avec unité~; on supposera pour simplifier que les degrés de $A$ sont positifs et que $A^0=\Z$. On désigne par $\Ahat$ le produit des $A^i$ (c'est un anneau commutatif avec unité), par $\Ahp$ le noyau de l'augmentation évidente $\Ahat\to A^0=\Z$. Alors $1+\Ahp$ est un sous-groupe du groupe multiplicatif des éléments inversibles de $\Ahat$. Considérons le groupe abélien
\begin{equation}\tag{1.14}
\widetilde A=\Z\times(1+\Ahp),
\end{equation}
dont les éléments seront notés $[n,x]$ avec $n\in\Z$ et $x=1+\sum_{i\geq 1}x^i\in1+\Ahp$ ($x^i\in A^i$). Nous écrirons ce groupe additivement, ainsi~:
\begin{equation}\tag{1.15}
[n,x]+[n',x']=[n+n',xx'].
\end{equation}
L'élément nul de ce groupe est $[0,1]$. Nous allons y introduire une structure de $\lambda$-anneau augmenté, compatible avec sa structure additive, l'augmentation étant $\varepsilon:[n,x]\longrightarrow n$. Le produit $[m,x][n,y]$ est complètement caractérisé par le fait qu'il s'écrive $[mn,x^n y^m(x*y)]$, où $x*y$ est un élément de $1+\Ahp$ dont les composantes s'expriment comme polynômes universels à coefficients entiers en les $x^i$ et les $y^i$~:
\begin{equation}\tag{1.16}
(x*y)^i=Q_i(x^1,\ldots,x^i;y^1,\ldots,y^i)\qquad (i\geq1)
\end{equation}
